%=====================================================================
\chapter{Design Science Research}\label{ch:designscience}
%=====================================================================
\locator{3}{Part~3 --- Research Designs for Computing}

\begin{outcomes}
By the end of this chapter you will be able to:
\begin{tight}
  \item \textbf{Distinguish} design science research from software development
    and from applied engineering.
  \item \textbf{Classify} an artefact as a construct, model, method or
    instantiation.
  \item \textbf{Apply} the DSRM process model to structure a project and a
    thesis.
  \item \textbf{Position} a contribution as routine design, exaptation,
    improvement or invention.
  \item \textbf{Select} an evaluation strategy using FEDS, and justify it
    against risk.
  \item \textbf{State} the design knowledge your artefact produces, separately
    from the artefact.
\end{tight}
\end{outcomes}

\begin{prereq}
\chref{contribution} --- especially the point that an artefact contributes by
carrying a claim, never by existing. \chref{philosophy} for pragmatism and for
why design science reasons abductively.
\end{prereq}

\begin{keyterms}
design science research $\bullet$ artefact $\bullet$ construct $\bullet$ model
$\bullet$ method $\bullet$ instantiation $\bullet$ design knowledge $\bullet$
relevance cycle $\bullet$ rigour cycle $\bullet$ design cycle $\bullet$
exaptation $\bullet$ formative and summative evaluation $\bullet$ naturalistic
and artificial evaluation $\bullet$ FEDS
\end{keyterms}

\section{The Tradition That Fits Computing Best, and Is Done Worst}

Most graduate computing projects build something. That makes design science
research the natural home for a great deal of our work --- and also makes it
the tradition most often invoked as cover for a project that is simply
development with a literature review attached.

\begin{definitionbox}[Design science research]
\term{Design science research} produces knowledge by building and evaluating an
artefact intended to solve a class of problem, and by articulating what the
building and evaluating taught that was not known before.

\medskip
Two clauses do the work. \emph{A class of problem}: not this hospital's data
pipeline, but the class of problems of which it is an instance. \emph{What it
taught}: the \term{design knowledge} --- a principle, a trade-off, a
demonstrated feasibility --- stated so that somebody solving a different
instance can use it without your code.
\end{definitionbox}

\begin{center}
\small
\begin{tabular}{@{}L{2.8cm}L{5.5cm}L{5.7cm}@{}}
\toprule
 & \textbf{Development / consultancy} & \textbf{Design science research} \\
\midrule
Target     & This client's problem
           & A class of problems this instance represents \\
Success    & The stakeholder is satisfied
           & A knowledge claim survives evaluation \\
Output     & A working system
           & An artefact \emph{and} transferable design knowledge \\
Novelty    & Not required
           & Required, and positioned against prior artefacts \\
Failure    & A failure
           & A result, if the design rationale was explicit \\
\bottomrule
\end{tabular}
\end{center}

\begin{alertbox}[The question that separates them]
Ask: \emph{if my artefact were deleted tomorrow, what would remain?}

\medskip
If the answer is ``nothing'', it was development. If it is ``the finding that
this class of problem yields to this principle, and the evidence for it'', it
was research. Write that sentence into your synopsis. It is what an examiner is
looking for and what most design-oriented theses never quite say.
\end{alertbox}

\section{What Counts as an Artefact}

March and Smith's four categories remain the clearest
taxonomy~\cite{march1995design}, and naming yours precisely tells you what to
evaluate.

\begin{center}
\small
\begin{tabular}{@{}L{2.4cm}L{4.6cm}L{7.0cm}@{}}
\toprule
\textbf{Type} & \textbf{What it is} & \textbf{Example, and what evaluation asks} \\
\midrule
Construct
  & Vocabulary and concepts
  & A taxonomy of failure modes in federated learning. Does it carve the space
    usefully? Can two people classify the same case identically? \\
Model
  & A representation of a problem and solution space
  & A formal model of trust propagation. Does it capture what matters, and does
    it predict? \\
Method
  & A process or algorithm
  & A procedure for partitioning clinical data across sites. Does it work, and
    is it better than the alternative on a stated dimension? \\
Instantiation
  & A working system embodying the above
  & The deployed framework. Does it demonstrate feasibility in a realistic
    setting? \\
\bottomrule
\end{tabular}
\end{center}

\begin{conceptbox}[Most theses claim an instantiation and contribute a method]
A common and easily fixed problem. The thesis presents a system, evaluates the
system, and buries the actual contribution --- a novel scheduling method, a new
way of decomposing the problem --- inside the implementation chapter.

\medskip
Separate them explicitly. The method is the contribution; the instantiation is
the evidence that the method is realisable. Written that way, the work
generalises, and a reader who will never run your code can still use your
result.
\end{conceptbox}

\section{Hevner's Guidelines, Read as Questions}

The seven guidelines~\cite{hevner2004design} are the most cited framing in the
field. They are more useful as a checklist of questions to answer in the
synopsis than as a description of process.

\begin{center}
\small
\begin{tabular}{@{}L{0.6cm}L{3.8cm}L{9.4cm}@{}}
\toprule
1 & Design as an artefact & What exactly is produced, and which of the four
                            types is it? \\
2 & Problem relevance     & Whose problem, and what does it cost them now? \\
3 & Design evaluation     & By what method, against what criteria, and how do
                            those criteria follow from the problem? \\
4 & Research contributions& What is now known that was not? State it as a
                            sentence, not as a system \\
5 & Research rigour       & Which existing theory and methods were drawn on ---
                            the rigour cycle \\
6 & Design as a search    & What alternatives were considered and rejected, and
                            on what grounds? \\
7 & Communication         & Can both a technical and a managerial audience take
                            something from it? \\
\bottomrule
\end{tabular}
\end{center}

\begin{pitfall}[Guideline 6 is the one that is skipped, and it is the one that matters]
Design is a search through a space of possible solutions. A thesis that
presents one design, fully formed, with no account of what else was considered,
has hidden the research. The alternatives you rejected --- and \emph{why} ---
are design knowledge, and often the most transferable part of the work.

\medskip
Keep a design log from day one: each significant decision, the options, the
criterion that decided it, and what you would need to have known to decide
differently. It costs minutes a week and becomes a thesis chapter.
\end{pitfall}

\section{The DSRM Process}

Peffers and colleagues give the process model most theses are structured
around~\cite{peffers2007design}.

\armfig{
\begin{tikzpicture}[
  st/.style={rounded corners=3pt, draw=primary!60, fill=primary!7,
             text width=2.3cm, align=center, font=\scriptsize,
             minimum height=1.5cm, inner sep=3pt},
  ar/.style={-{Stealth[length=2.2mm]}, draw=secondary, semithick}]

  \node[st] (p1) at (0,0)    {\textbf{1. Identify problem}\\ and motivate};
  \node[st] (p2) at (2.75,0) {\textbf{2. Define objectives}\\ of a solution};
  \node[st] (p3) at (5.5,0)  {\textbf{3. Design and develop}\\ the artefact};
  \node[st] (p4) at (8.25,0) {\textbf{4. Demonstrate}\\ in a context};
  \node[st] (p5) at (11.0,0) {\textbf{5. Evaluate}\\ against objectives};
  \node[st] (p6) at (13.75,0){\textbf{6. Communicate}};

  \foreach \a/\b in {p1/p2, p2/p3, p3/p4, p4/p5, p5/p6} {\draw[ar] (\a) -- (\b);}

  \draw[ar, dashed, draw=accent] (p5.south) .. controls +(0,-1.0) and +(0,-1.0) ..
       node[below, font=\scriptsize, text=accent] {iterate: back to design, or to objectives} (p3.south);
  \draw[ar, dashed, draw=accent] (p5.south) .. controls +(0,-1.7) and +(0,-1.7) .. (p2.south);

  \node[align=center, font=\scriptsize\itshape, text=neutral]
       at (6.9,1.5) {Entry point varies: problem-centred, objective-centred,
                     design-centred, or from an observed solution in use};
\end{tikzpicture}}
{The DSRM process. The iteration arrows are the substance: a design science
project that ran straight through once, with no loop back from evaluation, has
either been extraordinarily lucky or has not reported what actually
happened.}{dsrm}

\begin{conceptbox}[Report the iterations, do not hide them]
Theses routinely present DSRM as though it ran once, cleanly. Examiners know it
did not. The loops are where the design knowledge is: version 1 failed because
of X, which revealed that assumption Y was wrong, and version 2 addressed it.
That narrative is more convincing, more useful to a reader, and easier to write
than a retrospective fiction of a design that was right first time.
\end{conceptbox}

\section{Positioning the Contribution}

Gregor and Hevner's two-by-two~\cite{gregor2013positioning} settles the question
examiners ask: is this novel enough? Position it on solution maturity against
problem maturity.

\begin{center}
\small
\begin{tabular}{@{}L{2.6cm}L{5.6cm}L{5.6cm}@{}}
\toprule
 & \textbf{Problem is new} & \textbf{Problem is known} \\
\midrule
\textbf{Solution is new}
  & \emph{Invention.} Rare; a genuinely novel solution to a genuinely new
    problem
  & \emph{Improvement.} The most common good thesis: a better solution to a
    known problem. Requires comparison against the current best \\
\textbf{Solution exists}
  & \emph{Exaptation.} Extending a known solution into a new problem domain.
    Legitimate research \emph{if} the transfer was non-obvious and something
    was learned in making it work
  & \emph{Routine design.} Not a research contribution. Competent professional
    work \\
\bottomrule
\end{tabular}
\end{center}

\begin{alertbox}[Exaptation is where most ``applied'' theses live, and where they are won or lost]
``Applying deep learning to crop disease detection in Punjab'' is exaptation.
Whether it is research depends entirely on one thing: \textbf{was the transfer
non-obvious, and did making it work require solving something?}

\medskip
If the answer is that you fine-tuned a standard architecture on a local
dataset and it worked as expected, that is routine design, and no amount of
DSR vocabulary changes it (\chref{contribution}, non-contribution 1).

\medskip
If the answer is that local field conditions --- class imbalance across
smallholdings, imagery from low-cost handsets, labels from non-expert farmers
--- broke the standard approach in a specific way, and your artefact addresses
that, then the exaptation is genuine. The contribution is the characterisation
of \emph{why} it broke plus the design response. Say that in the abstract.
\end{alertbox}

\section{Evaluation: FEDS}

Evaluation is where design science theses most often fail, because
``demonstration'' is mistaken for evaluation. Showing that your system runs on
one example establishes existence, not merit.

The FEDS framework~\cite{venable2016feds} organises the choice along two
dimensions.

\armfigh{
\begin{tikzpicture}[
  qd/.style={draw=primary!50, rounded corners=2pt, text width=4.5cm,
             align=left, font=\scriptsize, inner sep=5pt, minimum height=2.3cm},
  ax/.style={font=\scriptsize\bfseries, text=secondary}]

  \node[qd, fill=primary!5]      (tl) at (0,0)     {\textbf{Formative + Artificial}\\[2pt]
        Early, cheap, controlled. Simulations, benchmarks, expert walkthroughs.
        Cheap to iterate; weak evidence of real-world merit};
  \node[qd, fill=secondary!7]    (tr) at (5.2,0)   {\textbf{Formative + Naturalistic}\\[2pt]
        Early, with real users on real tasks. Pilot deployments, action
        research. Expensive, but surfaces the assumptions that matter};
  \node[qd, fill=secondary!7]    (bl) at (0,-2.7)  {\textbf{Summative + Artificial}\\[2pt]
        Final, controlled. Controlled experiment against a baseline
        (\chref{experiments}); benchmark suite (\chref{benchmarking})};
  \node[qd, fill=goldhl!10]      (br) at (5.2,-2.7){\textbf{Summative + Naturalistic}\\[2pt]
        Final, in the field. Case study of real use (\chref{casestudy}).
        The strongest evidence of utility, and the hardest to obtain};

  \node[ax, rotate=90] at (-2.9,-1.35) {FORMATIVE $\rightarrow$ SUMMATIVE};
  \node[ax] at (2.6,1.55) {ARTIFICIAL $\rightarrow$ NATURALISTIC};
\end{tikzpicture}}
{The FEDS evaluation space. A thesis normally moves diagonally: cheap
artificial formative evaluation early to fix the design, then the most
naturalistic summative evaluation the constraints allow.}{feds}

\begin{center}
\small
\begin{tabular}{@{}L{3.6cm}L{10.7cm}@{}}
\toprule
\textbf{Strategy} & \textbf{Choose it when} \\
\midrule
Human risk and effectiveness
  & The main uncertainty is whether people will actually use it or benefit.
    Move to naturalistic evaluation early, even if crude \\
Technical risk and efficacy
  & The main uncertainty is whether the technique works at all. Stay
    artificial and rigorous for longer; naturalistic evaluation late \\
Quick and simple
  & Low risk, small artefact, limited resources. One or two evaluation
    episodes, honestly limited \\
Purely technical
  & There are no human users --- an algorithm, a protocol. Artificial
    evaluation throughout is appropriate, not a compromise \\
\bottomrule
\end{tabular}
\end{center}

\begin{conceptbox}[Choose the strategy from the risk, and say so]
The reason FEDS is worth using is that it makes evaluation choice
\emph{arguable}. A thesis that says ``we evaluated by simulation'' invites the
objection that simulation is unrealistic. A thesis that says ``the dominant
uncertainty in this design was technical efficacy rather than user acceptance,
so we followed a technical risk strategy: three rounds of artificial formative
evaluation, then a summative benchmark against the current best method, with a
single naturalistic case study to check external validity'' has pre-empted the
objection by answering it.
\end{conceptbox}

\begin{pitfall}[Four evaluations that are not evaluations]
\begin{tight}
  \item \textbf{Demonstration.} ``The system successfully processed the
    dataset.'' Establishes that it runs.
  \item \textbf{Self-comparison.} Comparing version 2 against version 1 of your
    own artefact. Shows you improved it; says nothing about the state of the
    art.
  \item \textbf{The satisfied-user survey.} Five students rate the interface
    4.2 out of 5. Who selected them, what were they comparing against, and what
    would a score of 2 have caused you to conclude?
  \item \textbf{Criteria chosen after the results.} If the evaluation criteria
    were not written into the protocol before the artefact was built, they
    were selected to be passed.
\end{tight}
\end{pitfall}

\begin{regbox}[What the regulations ask of a built artefact --- PhD \S14]
\S14(i) requires an ``original and innovative contribution to knowledge that
contributes to solving socioeconomic problems'' --- which is, almost exactly,
Hevner's problem relevance guideline.

\medskip
\S14.3(C), appropriateness of methods, requires that a dissertation ``justify
the use of methods and techniques to achieve study objectives'' and ``show
evaluation of obtained results in relation with study objectives''. For design
science that is a direct instruction: your evaluation criteria must be
traceable to your DSRM step 2 objectives.

\medskip
\S14.3(D) then asks for relevance to policy and practice, including that ``the
study output should be significant enough to be published or patented''.
\end{regbox}

\begin{traditions}{building an artefact}
\tradEMP{Treats the artefact as a treatment: what is the effect of using it,
compared with the alternative, measured how? This is the summative-artificial
quadrant, and \chref{experiments} supplies the machinery.}
\tradDSR{The artefact is the vehicle of the contribution, and design knowledge
is the contribution. Rigour comes from the design rationale and the evaluation,
not from statistics alone.}
\tradQUAL{Asks how the artefact is understood and appropriated by the people
using it --- often revealing that it is used for something other than its
designed purpose, which is itself a finding.}
\tradFORM{Asks what can be proved about the artefact: correctness, complexity,
guarantees under stated assumptions. Complements empirical evaluation rather
than replacing it.}
\end{traditions}

\begin{lab}[Position and plan an evaluation]
For an artefact you intend to build --- or one from a paper you admire:

\begin{tightnum}
  \item Name the artefact type: construct, model, method or instantiation. If
    you name two, say which is the contribution.
  \item Place it in the Gregor--Hevner two-by-two, and write one sentence
    defending the placement against the next cell down.
  \item Write the design knowledge sentence: ``if the artefact were deleted,
    what remains is\ldots''
  \item Identify the dominant risk --- technical or human --- and choose a FEDS
    strategy.
  \item Write the evaluation criteria \emph{now}, before building, and trace
    each to a DSRM step 2 objective.
\end{tightnum}

\medskip
Step 5 is the one that must be done in this order. Feed it into portfolio piece
P3.
\end{lab}

\begin{checkpoint}
\begin{tightnum}
  \item A scholar builds a recommender for a university library and reports
    that librarians liked it. Which FEDS quadrant is this, and what is the
    strongest objection?
  \item Distinguish exaptation from routine design with an example from your
    own area where the boundary is genuinely unclear.
  \item Why must evaluation criteria be written before the artefact is built?
    Name the specific bias.
  \item Your artefact is a method, and you built a system to test it. Which is
    the contribution, and how should the thesis chapters be ordered to make
    that clear?
\end{tightnum}
\end{checkpoint}

\begin{chaptersummary}
\begin{tight}
  \item Design science produces knowledge by building and evaluating an
    artefact for a \emph{class} of problem, and stating what was learned.
  \item If the artefact were deleted, what remains is the contribution. Write
    that sentence.
  \item Four artefact types; most theses claim an instantiation and actually
    contribute a method. Separate them.
  \item Hevner's guideline 6 --- design as search --- is the most skipped and
    the most transferable. Keep a design log.
  \item DSRM's iteration loops are the substance. Report them.
  \item Position with Gregor and Hevner: invention, improvement, exaptation,
    routine design. Applied theses are usually exaptation, and stand or fall on
    whether the transfer was non-obvious.
  \item Choose evaluation from risk using FEDS, and write the criteria before
    building. Demonstration, self-comparison and satisfied-user surveys are not
    evaluations.
\end{tight}
\end{chaptersummary}

\begin{reviewq}
  \item Design science is sometimes accused of being unfalsifiable: an artefact
    that fails is redesigned rather than refuted. Is this a fair charge?
    Construct the criticism at its strongest and then answer it.
  \item Naturalistic evaluation gives stronger evidence of utility but weaker
    control. For a specific artefact in your area, describe an evaluation
    programme that obtains both, and state its cost.
  \item Gregor and Hevner treat exaptation as a legitimate contribution.
    Where exactly is the line between exaptation and routine design, and who
    decides --- the author, the reviewer, or the problem?
  \item \reg{PhD Regs 2024}{\S14.3(D)} requires that research output be
    ``significant enough to be published or patented''. What does this imply
    for a design science thesis whose artefact is released as open source, and
    is the requirement well-suited to such work?
\end{reviewq}
